\chapter{Research and Implementation Roadmap}
\label{ch:roadmap}

The project will proceed through small executable examples rather than attempting a
complete logic and compiler at once.

\section{Milestones}

\begin{enumerate}
  \item Stabilise the event envelope, including identity, predecessor links, actor
        endpoints, kind, payload, and a disciplined account of time.
  \item Give the grammar algebra a trace semantics and establish laws for sequence,
        choice, guarded context, and interleaving.
  \item Define parameter binding and participant projection, then test both on
        authentication and ordinary web request protocols.
  \item Preserve named scenario boundaries while reading their references into one
        grammar; render interaction diagrams and residual state machines per scenario.
  \item Treat independent requirements as trace-language intersections and report
        minimal inconsistent requirement sets.
  \item Generate finite transition systems and evaluate the CTL fragment already
        represented by the LeanFM library.
  \item Compare past-trace predicates with explicit past operators.
  \item Extend the temporal core with separately typed knowledge and policy
        modalities. Begin with bounded Kerberos histories, explicit local
        observations and nonempty ideal futures, and checked counterexamples;
        do not present the narrated example as an implemented checker.
  \item Render declarative pie and line charts as reductions over observable events.
  \item Connect specifications to a Lean web service and generated protocol
        schemas, preserving one source of typed domain truth.
\end{enumerate}

\section*{Self-hosting consolidation}

The canonical requirements entry point now composes the existing LeanFM ChatUI
specification with a proposed revision-bound requirements-to-PR workflow. The
Worker protocol remains an example, not a specification of the whole product.
\texttt{make selfhost-check} validates the typed composition and regenerates
\texttt{SELF\_HOSTING\_REVIEW.md}, including its interaction diagram and state
machine. A component registry and fresh-session handoff preserve decisions
without requiring conversation history.

The consolidation records demand-driven modal indexes, optional probability
with explicit interpretation, checker outcomes, acceptance invalidation, and
PR evidence. These are pending requirements, not implemented checker features
or proven predicates. Error protocols, executable formula semantics, and
complete implementation traceability still require refinement and both
acceptance gates before claiming a complete one-shot specification.

\section{Proof obligations}

The initial metatheory should make the following questions precise:

\begin{itemize}
  \item Does a projected local behaviour faithfully implement its part of the
        global grammar?
  \item Under what finiteness restrictions is CTL checking decidable and executable?
  \item Which guards can be compiled into finite monitor state?
  \item When does modular composition preserve an established temporal property?
  \item Which fairness and timing assumptions are required by a liveness theorem?
  \item Which observation and policy relations justify each knowledge,
        permission, and obligation assertion, and which bounds make their
        combination with CTL executable?
\end{itemize}

The intended outcome is not merely a verifier. It is a development loop in which
writing the behaviour is the fastest way to obtain working, reviewable software.
